\section{Implementación de la Solución}

Este capítulo documenta cómo se ejecutó el diseño del capítulo anterior. Toda la infraestructura está codificada en Terraform, toda la configuración en manifiestos YAML versionados en Git, y toda prueba está automatizada en CronJobs de Kubernetes ---condición necesaria para que los resultados sean reproducibles y atribuibles al CNI \cite{ref7}.

\subsection{Aprovisionamiento de Infraestructura con Terraform}

El aprovisionamiento de servidores y redes se realizó de forma completamente automatizada mediante Terraform (IaC). Al describir la infraestructura en archivos versionados, Terraform la recrea de manera idéntica en cada ejecución y elimina la posibilidad de que variaciones manuales en el aprovisionamiento introduzcan diferencias entre los entornos de cada CNI \cite{ref7}.

\subsubsection{Estructura del Código Terraform (K8s-bootstrap/00.bootstrap)}

% [EDIT] Tabla booktabs y caption superior.
\begin{table*}[htbp]
\centering
\scriptsize
\renewcommand{\arraystretch}{1.25}
\setlength{\tabcolsep}{4pt}
\caption{Archivos Terraform del módulo de aprovisionamiento y su función específica.}
\label{tab:terraform-aprovisionamiento}
\begin{tabularx}{\textwidth}{|l|Y|}
\hline
\textbf{Archivo Terraform} & \textbf{Responsabilidad en la Infraestructura} \\
\hline
\texttt{main.tf} & Define los recursos principales: Droplets de DigitalOcean y red VPC privada. \\
\hline
\texttt{variables.tf} & Centraliza todos los parámetros configurables: número de nodos (\texttt{agent\_count = 2}), tamaño de instancias, región y proveedor CNI. \\
\hline
\texttt{ssh\_keys.tf} & Automatiza la gestión de llaves SSH para el acceso seguro entre nodos. \\
\hline
\texttt{firewall.tf} & Configura las reglas de firewall a nivel de nube: primera línea de seguridad perimetral. \\
\hline
\texttt{outputs.tf} & Genera automáticamente el archivo kubeconfig con las credenciales del clúster K3s. \\
\hline
\end{tabularx}
\end{table*}

El proceso de aprovisionamiento tiene tres etapas: \texttt{terraform plan} (compara estado actual vs. código), \texttt{terraform apply} (crea Droplets, VPC y firewall) y \texttt{terraform output} (extrae IPs y kubeconfig). Terraform mantiene un archivo \texttt{terraform.tfstate} que permite destruir y recrear toda la infraestructura de manera controlada, garantizando la repetibilidad científica: cada CNI se evaluó en condiciones idénticas de arranque.

\subsection{Configuración del Clúster K3s en Alta Disponibilidad}

Una vez aprovisionados los servidores, se instaló y configuró K3s, una distribución liviana de Kubernetes, con topología de 1 nodo Master + 2 nodos Worker dedicados.

\subsubsection{Base de Datos Externa para el Plano de Control}

El plano de control se configuró con una base de datos PostgreSQL gestionada de DigitalOcean (también provisionada via Terraform) como datastore externo. Esto elimina la complejidad de mantener un clúster etcd y garantiza un plano de control sin punto único de fallo (SPOF).

\subsubsection{Instalación: clúster sin CNI para plugins externos y modo nativo para Flannel}

La implementación establece una distinción metodológica deliberada: Flannel se evalúa en su modo de despliegue más habitual ---integrado nativamente en K3s con \texttt{flannel-iface: eth1}---, mientras que Calico, Cilium y Antrea se instalan como plugins externos sobre un clúster sin CNI. La asimetría refleja las condiciones reales de uso: Flannel es el CNI predeterminado de K3s y la mayoría de equipos lo adopta sin configuración adicional, en tanto que los otros tres requieren instalación explícita y declarativa. Evaluar cada CNI en su escenario de despliegue típico aumenta la validez ecológica del experimento; debe considerarse, no obstante, que Flannel se beneficia de la integración nativa de K3s ---gestión automática de CIDR y \emph{hooks} de kubelet--- que los plugins externos deben replicar por su cuenta.

Para los CNIs externos, la configuración crítica es el parámetro \texttt{--flannel-backend=none}, que desactiva completamente el CNI integrado de K3s y evita interferencias con el plugin bajo evaluación. Adicionalmente, se configura \texttt{--disable-network-policy} para que K3s no intente implementar NetworkPolicies por sí mismo.

La configuración resultante en \texttt{/etc/rancher/k3s/config.yaml} para CNIs externos incluye:
\begin{lstlisting}
flannel-backend: none
disable-network-policy: true
\end{lstlisting}

% [EDIT] Tabla booktabs y label consistente.
\begin{table*}[htbp]
\centering
\scriptsize
\renewcommand{\arraystretch}{1.25}
\setlength{\tabcolsep}{4pt}
\caption{Instalación de CNIs: método y configuración clave en \texttt{user\_data/ks3\_server\_init.sh}.}
\label{tab:cni-instalacion}
\begin{tabularx}{\textwidth}{|l|p{3.2cm}|Y|}
\hline
\textbf{CNI} & \textbf{Método de instalación} & \textbf{Configuración clave} \\
\hline
Flannel & Nativo K3s & \texttt{flannel-iface: eth1} (interfaz VPC de DigitalOcean) \\
\hline
Calico & Tigera Operator v3.29.1 & \texttt{encapsulation: VXLAN}, CIDR \texttt{10.42.0.0/16}, \texttt{nodeAddressAutodetectionV4.interface: eth1} \\
\hline
Cilium & Helm 1.16.5 & \texttt{routingMode=tunnel}, \texttt{tunnelProtocol=vxlan}, \texttt{ipam.mode=kubernetes}, \texttt{devices=eth1}, \texttt{operator.replicas=1} \\
\hline
Antrea & Manifest v2.2.0 & \texttt{transportInterface: eth1}, \texttt{serviceCIDR: 10.43.0.0/16} \\
\hline
\end{tabularx}
\end{table*}

El protocolo de rotación de CNIs tuvo cuatro pasos: (1) \texttt{kubectl delete} para eliminar recursos del CNI actual; (2) verificación de que no quedan interfaces de red anteriores en los nodos; (3) reinicio de los nodos Worker para limpiar estado en memoria; (4) aplicación del manifiesto del nuevo CNI via ArgoCD. Este protocolo garantiza que cada CNI inicia sobre un nodo completamente limpio, sin residuos de configuraciones anteriores que pudieran contaminar las métricas del siguiente ciclo de pruebas.

\subsection{Orquestación y Entrega Continua: GitOps con ArgoCD}

Con el clúster K3s funcionando y sin CNI, el siguiente componente instalado fue ArgoCD. ArgoCD implementa el paradigma GitOps: el repositorio Git es la única fuente de verdad sobre el estado del clúster. Todo lo que debe existir en el clúster está definido en archivos YAML dentro del repositorio \texttt{K8s-Tesis-Apps}. Esta centralización en Git facilita la auditoría del experimento: cualquier cambio en la configuración queda registrado con autor, fecha y motivo, haciendo el proceso de prueba completamente trazable.

\subsubsection{El Patrón App of Apps}

La implementación de ArgoCD utiliza el patrón App of Apps: existe UNA sola aplicación raíz (ubicada en \texttt{K8s-bootstrap/30.apps-of-apps-install}) que ArgoCD conoce directamente, y esa aplicación raíz contiene las definiciones de todas las demás aplicaciones (CNIs, stack de métricas, benchmarks, políticas de red).

\subsubsection{Configuración de selfHeal y prune}

\textbf{\emph{selfHeal: true}} --- Ante la modificación manual de un recurso del clúster ---por ejemplo, la eliminación de una NetworkPolicy---, ArgoCD detecta la diferencia con el repositorio Git y restaura el estado original, lo que resultó crítico para la integridad del experimento.

\textbf{\emph{prune: true}} --- Al eliminar un archivo YAML del repositorio Git, ArgoCD elimina también el recurso correspondiente del clúster, lo que evita la acumulación de recursos huérfanos.

\begin{lstlisting}
# Definicion de la Application raiz en ArgoCD (App of Apps)
apiVersion: argoproj.io/v1alpha1
kind: Application
metadata:
  name: apps-of-apps
  namespace: argocd
spec:
  source:
    repoURL: https://github.com/equipo/K8s-Tesis-Apps
    targetRevision: main
    path: apps/
  syncPolicy:
    automated:
      selfHeal: true   # Restaura cambios manuales automaticamente
      prune: true      # Elimina recursos borrados del repo
\end{lstlisting}

\subsubsection{Overlays de Kustomize para los CNIs}

Cada CNI tiene una configuración ligeramente diferente. Para gestionarlas sin duplicar código, se utilizó Kustomize con el patrón base/overlay: existe una configuración base con los parámetros comunes, y sobre esa base se aplican overlays específicas para cada CNI.

Las apps gestionadas por ArgoCD incluyen: \texttt{cni-test-iperf} (benchmarks de throughput y latencia), \texttt{cni-resource-exporter} (exportación de recursos desde Prometheus), \texttt{cni-np-test} (pruebas de Network Policies, creada solo cuando \texttt{cni\_provider != "flannel"} y el caso de uso está definido), y \texttt{metrics} (Prometheus + Grafana).

\subsection{Implementación del Stack de Observabilidad}

El stack de observabilidad se implementó mediante el \texttt{kube-prometheus-stack} (Helm chart) que instala Prometheus, Grafana y exportadores preconfigurados para Kubernetes.

\subsubsection{Prometheus: El Recolector de Métricas}

Prometheus opera con un modelo pull: visita periódicamente los endpoints de métricas (cada 15 segundos en la configuración del proyecto). Para que Prometheus sepa qué métricas recolectar de los agentes CNI, se implementaron \texttt{ServiceMonitors} (CRD), apuntando específicamente al endpoint de métricas de cada agente. Un intervalo de 15 segundos proporciona suficiente granularidad para detectar picos de consumo durante las ráfagas de tráfico iperf3, sin generar una carga de scraping que interfiera con las métricas bajo observación.

Las métricas clave recolectadas del agente CNI incluyen:
\begin{itemize}
\item \path|container_cpu_usage_seconds_total| (filtrado por pod del agente): costo computacional del CNI (criterio C2 del marco MCDA).
\item \path|container_memory_rss| (filtrado por pod del agente): memoria física real del proceso del agente CNI.
\item \path|node_network_transmit_bytes_total| / \path|node_network_receive_bytes_total|: bytes totales de red por nodo.
\item \path|kube_pod_info| + \path|kube_pod_status_phase|: estado del agente CNI para detectar reinicios durante las pruebas.
\end{itemize}

\subsubsection{Grafana: Los Dashboards de Red}

% latex-paper-en · deai (Line 142) [Severity: P0] [Priority: P0]: lista paralela mecánica
Grafana centraliza la lectura operativa de las series de Prometheus. Sus paneles agrupan throughput, latencia, consumo del agente CNI y score MCDA. Esta correlación vincula ráfagas iperf3 con CPU, RSS y el ranking calculado por \texttt{procesador.js}.

\subsubsection{Resource Usage Exporter: Extracción Automatizada de Datos}

% latex-paper-en · deai (Line 147) [Severity: P1] [Priority: P1]: tono operativo repetitivo
El \texttt{cni-resource-usage-exporter} corre como CronJob en \path|K8s-Tesis-Apps/grafana-exporter|. ArgoCD sincroniza el recurso mediante la app \texttt{cni-resource-exporter}. El CronJob:
\begin{itemize}
\item Ejecuta la extracción en los minutos 25 y 55 de cada hora, tras cada ventana de benchmarks.
\item Consulta directamente la API HTTP de Prometheus (\texttt{/api/v1/query\_range}) --- no interactúa con Grafana.
\item Recolecta CPU\% por nodo, RAM\% por nodo, RAM usada, CPU cores y RAM del agente CNI.
\item Genera \texttt{raw/}, \texttt{csv/}, \texttt{images/} y \texttt{cni\_resource\_summary.json} por corrida.
\item Persiste cada salida con timestamp en \path|K8S-CNI-Results/results/cni-benchmarks/{cni}/resource_usage_nodes/|.
\end{itemize}

\subsection{Implementación del Banco de Pruebas de QoS}

\subsubsection{Escenario de Rendimiento con iperf3}

% latex-paper-en · deai (Line 161) [Severity: P1] [Priority: P1]: frase genérica
El banco de throughput se define en \texttt{K8s-Tesis-Apps/cni\_test\_iperf}:

\textbf{El iperf3-server (Deployment):} Un Deployment ejecuta una réplica de iperf3 en modo servidor. Un Service ClusterIP expone un nombre DNS estable para el cliente.

\textbf{El iperf3-client (CronJob):} La programación \texttt{0,30 * * * *} activa el cliente en los minutos 0 y 30. \texttt{concurrencyPolicy: Forbid} bloquea corridas simultáneas. Cada activación crea un Pod y genera tráfico TCP durante 300 segundos. La salida conserva formato JSON con \texttt{iperf3 -t 300 -i 10 -J}. La retención máxima queda en 5 corridas (\texttt{RETENTION\_MAX\_RUNS="5"}).

La anti-afinidad del cliente frente al servidor se implementó como \texttt{preferredDuringSchedulingIgnoredDuringExecution} (suave, weight 100), forzando tráfico inter-nodo en la práctica. Esta ubicación garantiza que todo el tráfico de benchmark atraviese el mecanismo de encapsulación del CNI entre nodos, que es exactamente lo que se quiere medir.

\begin{lstlisting}
# iperf-client-cronjob.yaml -- fragmento relevante
schedule: '0,30 * * * *'   # Minutos 0 y 30
concurrencyPolicy: Forbid
affinity:
  podAntiAffinity:
    preferredDuringSchedulingIgnoredDuringExecution:
    - weight: 100
      podAffinityTerm:
        labelSelector:
          matchLabels:
            app: iperf-server
        topologyKey: kubernetes.io/hostname
# iperf3 -c iperf-server-svc -t 300 -i 10 -J
\end{lstlisting}

\subsubsection{Cliente de Latencia Personalizado}

% latex-paper-en · deai (Line 188) [Severity: P1] [Priority: P1]: frase temporal mecánica
El CronJob \texttt{latency-client-cronjob.yaml} ejecuta la prueba de latencia. Su programación corre en los minutos 10 y 40 de cada hora. Esa separación evita solapamiento con el inicio del throughput.

El cliente registra TCP Connect con \texttt{nc -z -w 2} y timestamp en ms. Cada corrida toma 30 muestras y calcula min/avg/max/mdev. El JSON normalizado conserva campos de latencia, muestras intentadas y muestras exitosas.

% [EDIT] Tabla booktabs y label consistente.
\begin{table*}[htbp]
\centering
\scriptsize
\renewcommand{\arraystretch}{1.25}
\setlength{\tabcolsep}{4pt}
\caption{Métricas de QoS capturadas por el banco de pruebas y su relevancia telemática.}
\label{tab:metricas-qos-implementacion}
\begin{tabularx}{\textwidth}{|p{4cm}|p{2.4cm}|Y|}
\hline
\textbf{Métrica Capturada} & \textbf{Unidad} & \textbf{Relevancia para QoS Telemático} \\
\hline
Throughput TCP sostenido & Gbps / Mbps & Capacidad efectiva de la red. \\
\hline
Latencia mínima TCP connect & ms & Overhead mínimo del CNI sin congestión. \\
\hline
Latencia promedio TCP connect & ms & Comportamiento típico de solicitudes. \\
\hline
Latencia máxima / Jitter & ms & Peor caso y variabilidad. Crítico para tiempo real. \\
\hline
Retransmisiones TCP & count / sesión & Indicador de problemas en encapsulamiento o congestión. \\
\hline
CPU del agente CNI & millicores (m) & Impacto directo en costo de infraestructura cloud. \\
\hline
RAM del agente CNI (RSS) & MiB & Huella de memoria del agente CNI por nodo. \\
\hline
\end{tabularx}
\end{table*}

\subsection{Procesamiento Estadístico y Aplicación del Modelo MCDA}

\subsubsection{Limpieza con el Rango Intercuartílico (IQR)}

% latex-paper-en · deai (Line 226) [Severity: P2] [Priority: P2]: secuencia formulaica
\texttt{docs/procesador.js} transforma datos crudos en datos consolidados. El script descubre los CNIs disponibles bajo \texttt{results/cni-benchmarks}. Luego extrae métricas desde JSON y CSV. Finalmente aplica IQR:

\begin{lstlisting}
// procesador.js -- Funcion de limpieza estadistica IQR
function calcularEstadisticasLimpias(valores) {
  const sorted = [...valores].sort((a, b) => a - b);
  const q1 = sorted[Math.floor(sorted.length * 0.25)];
  const q3 = sorted[Math.floor(sorted.length * 0.75)];
  const iqr = q3 - q1;
  const lowerBound = q1 - 1.5 * iqr;
  const upperBound = q3 + 1.5 * iqr;
  const valoresLimpios = valores.filter(
    v => v >= lowerBound && v <= upperBound);
  const promedio = valoresLimpios.reduce(
    (sum, v) => sum + v, 0) / valoresLimpios.length;
  return { promedio,
    outliersPorcentaje: ((valores.length - valoresLimpios.length)
      / valores.length * 100).toFixed(1)
  };
}
\end{lstlisting}

% latex-paper-en · deai (Line 249) [Severity: P1] [Priority: P1]: causalidad débil
El pipeline descarta valores atípicos antes de calcular promedios. Dado que un evento cloud anómalo puede sesgar una corrida, el filtro estabiliza el puntaje MCDA del CNI.

\subsubsection{Normalización y Cálculo del Score MCDA}

% latex-paper-en · deai (Line 254) [Severity: P1] [Priority: P1]: estructura repetida
Después de la limpieza, el pipeline aplica normalización min-max. Cada valor se escala al rango $[0, 1]$. En métricas positivas, como throughput, la normalización es directa. En métricas negativas, como latencia o CPU, la escala se invierte.

Cuando la métrica penaliza, el valor más bajo recibe 1 y el más alto recibe 0. La escala común compara Gbps, milisegundos y millicores sin sesgo por magnitud numérica.

\begin{lstlisting}
// procesador.js -- Calculo del score MCDA por perfil
const perfiles = {
  fintech:   { throughput: 0.15, latencia: 0.20,
               retransmisiones: 0.15, cpu: 0.10,
               ram: 0.10, networkPoliciesL7: 0.30 },
  streaming: { throughput: 0.35, latencia: 0.35,
               retransmisiones: 0.20, cpu: 0.10,
               ram: 0.00, networkPoliciesL7: 0.00 },
  iot:       { throughput: 0.25, latencia: 0.20,
               retransmisiones: 0.15, cpu: 0.30,
               ram: 0.10, networkPoliciesL7: 0.00 }
};
\end{lstlisting}

% latex-paper-en · deai (Line 274) [Severity: P2] [Priority: P2]: transición sin carga lógica
El pipeline calcula overhead de seguridad cuando detecta \texttt{with\_network\_policy}. En ese caso, compara la métrica base contra la métrica con política activa.

\subsection{Prototipo Recomendador: Aplicación Web SPA}

\subsubsection{Tecnologías de Implementación}

La SPA fue implementada con React~19 (gestión de estado por componentes), Vite~7 (bundler; build de producción vía \texttt{build:pages} con \texttt{--outDir ../docs/recommender}) y Tailwind CSS~3 (estilos utilitarios directamente en JSX, sin CSS personalizado).

% latex-paper-en · deai (Line 305) [Severity: P2] [Priority: P2]: secuencia plana
\texttt{npm run sync:data} ejecuta \texttt{procesador.js} y genera \texttt{cni-data.json}. Después, \texttt{export-spa-data.cjs} inyecta los datos en el bundle. \texttt{recommendationModel.js} concentra la lógica de recomendación.

\subsubsection{Flujo de la Aplicación}

\begin{enumerate}
\item \textbf{Selección de Perfil:} La pantalla principal muestra tres tarjetas: Fintech/Seguridad, Streaming/Rendimiento e IoT/Recursos.
\item \textbf{Visualización de Resultados:} La aplicación carga \texttt{cni-data.json} y calcula el ranking MCDA para el perfil seleccionado. El CNI con mayor puntaje queda marcado como recomendación.
\item \textbf{Justificación Técnica:} La interfaz expone las métricas brutas promedio de cada CNI en gráficas comparativas.
\end{enumerate}

\subsection{Implementación de Network Policies y Seguridad Zero Trust}

\subsubsection{Políticas Default Deny}

% latex-paper-en · deai (Line 320) [Severity: P0] [Priority: P0]: paralelismo mecánico
Default Deny operacionaliza el modelo Zero Trust. La política bloquea todo el tráfico inicial y habilita solo comunicaciones necesarias. Con esta regla, toda comunicación no autorizada falla sin intervención del operador.

El siguiente manifiesto aplica la política a todos los Pods del namespace \texttt{production}:

\begin{lstlisting}
# default-deny-all.yaml
apiVersion: networking.k8s.io/v1
kind: NetworkPolicy
metadata:
  name: default-deny-all
  namespace: production
spec:
  podSelector: {}     # aplica a TODOS los Pods del namespace
  policyTypes:
  - Ingress           # Bloquea todo el trafico entrante
  - Egress            # Bloquea todo el trafico saliente
\end{lstlisting}

\subsubsection{Micro-segmentación por Capas de Aplicación}

% latex-paper-en · deai (Line 341) [Severity: P1] [Priority: P1]: frase genérica
Sobre Default Deny, las políticas específicas acotan el tráfico permitido. Las reglas \texttt{frontend $\rightarrow$ backend $\rightarrow$ database} usan label selectors.

En Cilium, una \texttt{CiliumNetworkPolicy} de capa 7 admite solo métodos \texttt{GET} y \texttt{HEAD} desde frontend hacia backend. En Antrea, una \texttt{ClusterNetworkPolicy} en el Tier \texttt{securityops} prioriza sobre cualquier \texttt{NetworkPolicy} de namespace.

Esta estratificación evalúa seguridad granular en condiciones reales. También cuantifica si esa granularidad introduce overhead detectable, uno de los criterios del modelo MCDA.

\subsubsection{Validación Automatizada de Políticas}

Se implementó un pipeline de validación que emplea \texttt{nc} (netcat) dentro de Pods de prueba para verificar la conectividad. Para cada par de Pods, el pipeline define una expectativa ---la conexión debe permitirse o debe bloquearse--- y reporta un error cuando el comportamiento observado no coincide con ella.

% [EDIT] Tabla booktabs y label para referencia.
\begin{table*}[htbp]
\centering
\scriptsize
\renewcommand{\arraystretch}{1.25}
\setlength{\tabcolsep}{4pt}
\caption{Casos de prueba del pipeline de validación de Network Policies.}
\label{tab:validacion-network-policies}
\begin{tabularx}{\textwidth}{|p{5cm}|Y|}
\hline
\textbf{Prueba de Validación} & \textbf{Resultado Esperado y Justificación} \\
\hline
frontend $\rightarrow$ backend (puerto 8080) & PERMITIDO --- Comunicación necesaria para el funcionamiento de la aplicación. \\
\hline
frontend $\rightarrow$ database (puerto 5432) & BLOQUEADO --- El frontend no debe acceder directamente a la base de datos. \\
\hline
backend $\rightarrow$ database (puerto 5432) & PERMITIDO --- El backend necesita leer y escribir datos. \\
\hline
database $\rightarrow$ cualquier Pod & BLOQUEADO --- La base de datos nunca debe iniciar conexiones salientes. \\
\hline
Pod externo $\rightarrow$ namespace production & BLOQUEADO --- Sin una regla explícita, ningún Pod externo puede entrar. \\
\hline
\end{tabularx}
\end{table*}

\subsection{Estructura de Repositorios y Entregables Técnicos}

% latex-paper-en · deai (Line 380) [Severity: P1] [Priority: P1]: repite contenido tabular
La Tabla~\ref{tab:repositorios-entregables} resume la organización del proyecto en tres repositorios Git.

% [EDIT] Tabla booktabs y label consistente.
\begin{table*}[htbp]
\centering
\scriptsize
\renewcommand{\arraystretch}{1.25}
\setlength{\tabcolsep}{4pt}
\caption{Estructura de los tres repositorios del proyecto y su relación con los objetivos específicos.}
\label{tab:repositorios-entregables}
\begin{tabularx}{\textwidth}{|p{3cm}|Y|p{5.2cm}|}
\hline
\textbf{Repositorio} & \textbf{Contenido Principal} & \textbf{Relación con los Objetivos} \\
\hline
\texttt{K8s-bootstrap} & Código Terraform (\texttt{00.bootstrap}), scripts de instalación K3s HA (\texttt{10.k3s-ha-do}), manifiestos de ArgoCD (\texttt{20.argo-cd-install}) y App of Apps (\texttt{30.apps-of-apps-install}). & OE1: Provee la infraestructura controlada necesaria para las pruebas. \\
\hline
\texttt{K8s-Tesis-Apps} & Overlays Kustomize para los 4 CNIs, stack de observabilidad (\texttt{apps/metrics}), benchmarks iperf3 (\texttt{cni\_test\_iperf}), cliente de latencia, Network Policies y Grafana Exporter. & OE1, OE2 y OE3: Implementa las pruebas, la seguridad y el sistema de recolección. \\
\hline
\texttt{K8S-CNI-Results} & Archivos JSON/CSV con resultados por corrida, \texttt{procesador.js} (modelo MCDA) y la SPA recomendadora (\texttt{cni-recommender-spa}). & OE3 y OE4: Implementa el modelo de comparación objetiva y el prototipo. \\
\hline
\end{tabularx}
\end{table*}

% latex-paper-en · deai (Line 404) [Severity: P2] [Priority: P2]: cadena densa
La arquitectura mantiene un flujo unidireccional. \texttt{K8s-bootstrap} crea infraestructura; \texttt{K8s-Tesis-Apps} despliega aplicaciones y genera métricas. \texttt{K8S-CNI-Results} almacena, procesa y visualiza resultados.

\subsection{Desafíos de Implementación y Soluciones Adoptadas}

% [EDIT] Tabla booktabs y label consistente.
\begin{table*}[htbp]
\centering
\scriptsize
\renewcommand{\arraystretch}{1.25}
\setlength{\tabcolsep}{4pt}
\caption{Desafíos técnicos encontrados durante la implementación y soluciones adoptadas.}
\label{tab:desafios-implementacion}
\begin{tabularx}{\textwidth}{|Y|Y|}
\hline
\textbf{Desafío Encontrado} & \textbf{Solución Técnica Implementada} \\
\hline
MTU mismatch al instalar Cilium: Pods no podían comunicarse, paquetes se fragmentaban silenciosamente en la VPC. Flannel y Calico detectan automáticamente la MTU de la interfaz subyacente en sus versiones actuales; Cilium requirió configuración explícita porque su dataplane eBPF no hereda el valor de MTU del sistema operativo sin intervención. & Se calculó la MTU efectiva: 1500 (VPC) $-$ 50 (overhead VXLAN) $=$ 1450. Se configuró el parámetro \texttt{mtu: 1450} en el overlay de Kustomize de Cilium. \\
\hline
ArgoCD en estado OutOfSync tras la rotación de CNI: Pods perdían conectividad durante el cambio. & Se implementó una secuencia de rotación segura: poner ArgoCD en modo Suspended, rotar el CNI, verificar conectividad, reactivar ArgoCD. \\
\hline
Outliers frecuentes en pruebas de madrugada: DigitalOcean mostraba throttling de CPU en horas de alta demanda del datacenter. & Se configuró una ventana de pruebas exclusiva para horario laboral. El algoritmo IQR filtra automáticamente estos outliers. \\
\hline
\texttt{procesador.js} necesitaba datos de Network Policies L7 que iperf3 no puede medir (capa 4). & Se añadió una variable binaria (soporta/no soporta) al modelo MCDA para la capacidad de Network Policies L7. Cilium recibe valor 1, los demás 0. \\
\hline
\end{tabularx}
\end{table*}

\subsection{Flujo completo de ejecución}

% latex-paper-en · deai (Line 434) [Severity: P1] [Priority: P1]: repite tabla cercana
El flujo completo sigue las etapas de la Tabla~\ref{tab:flujo-implementacion}. Terraform prepara la infraestructura; K3s y ArgoCD habilitan el plano de ejecución. Los CronJobs producen datos, \texttt{procesador.js} consolida resultados y la SPA presenta la recomendación.

